\chapter{巨正则系综}\label{ch:b3:grand-canonical}

湿润空气中的一滴水既不增大，也不缩小；
氧气在肺部与血红蛋白结合，在肌肉中又解离；
电池驱动电子流过手机，是因为电子
在两个电极中的“价格”不同。在这些情形中，
被平衡的并非温度，而是\emph{\omterm{thm:b3:grand-canonical:distribution}{化学势}}
$\mu$：多加入一个粒子所需的\omterm{prop:b3:canonical-ensemble:machine}{自由能}代价。
本章打开能量库的第二道阀门——粒子与能量一样参与交换——
并建立\omterm{thm:b3:grand-canonical:distribution}{巨正则系综}。
它的回报令人惊叹：费米子或玻色子单个量子能级的占据数，
只需两行推导即可得到，
而下一章的主要公式也随之而来。它的应用范围很广，
从防毒面具与催化转化器中的\omterm{prop:b3:grand-canonical:langmuir}{吸附}等温线，直到\omterm{prop:b3:grand-canonical:saha}{萨哈方程}；
1925 年，哈佛大学的一名博士生借助它读取恒星的温度，
并出人意料地发现，
宇宙主要由氢构成。

\section{粒子加入协商}

\begin{theorem}[巨正则分布]\label{thm:b3:grand-canonical:distribution}
\index{巨正则系综}\index{化学势}
与能量库交换能量\emph{与粒子}的体系，能量库温度为
$T$、\omterm{thm:b3:grand-canonical:distribution}{化学势}为 $\mu$，体系处于状态 $s$
（能量 $E_s$，粒子数 $N_s$）的概率为
\[
\mathcal P_s = \frac{\eu^{-\beta(E_s - \mu N_s)}}{\Xi} , \qquad
\Xi = \sum_s\eu^{-\beta(E_s - \mu N_s)} ,
\]
此即\emph{巨\omterm{thm:b3:canonical-ensemble:boltzmann}{配分函数}}。由 $\Xi$ 可得 $\langle N\rangle
= k_{\text{B}}T\,\partial\ln\Xi/\partial\mu$，并且\emph{巨
势} $\Phi = -k_{\text{B}}T\ln\Xi$ 满足 $\Phi = -PV$。
\end{theorem}

\begin{proof}[部分证明]
仿\cref{thm:b3:canonical-ensemble:boltzmann}，将
能量库的熵在取出能量与粒子两方面展开：
$S_{\text{能量库}}(E - E_s, N - N_s) \approx S_{\text{能量库}} - E_s/T +
\mu N_s/T$，利用 $\partial S/\partial N = -\mu/T$
（见\cref{thm:b3:microcanonical-ensemble:temperature}）。平均值
由求导得到；$\Phi = -PV$ 在本阶段
不加证明地承认（它由广延性可得）。
\end{proof}

\begin{definition}[$\mu$ 的含义]\label{def:b3:grand-canonical:mu}
\index{化学势!含义}
$\mu$ 是加入一个粒子所需的\omterm{prop:b3:canonical-ensemble:machine}{自由能}代价，
在 $T$ 与 $V$ 固定时。粒子自发地从高 $\mu$ 流向低
$\mu$——扩散、渗透、蒸发、化学反应与
电流，都是 $\mu$ 梯度的弛豫过程。平衡时
任意两相或两处之间，$\mu$ 在两侧相等——
这是（与 $T$、$P$ 一起）支配共存的第三个等式。对
一个经典理想气体，
\[
\mu = k_{\text{B}}T\,\ln\!\big(n\lambda_T^3\big) < 0
\quad\text{（稀薄气体）} :
\]
向稀薄气体中加入一个粒子会\emph{降低}\omterm{prop:b3:canonical-ensemble:machine}{自由能}，
因为熵增超过任何能量代价——这正是
驱动蒸发进入干燥空气的信号。电池则是一台 $\mu$ 机器：
其电压为两电极之间每个电子的\omterm{thm:b3:grand-canonical:distribution}{化学势}降，
$U = \Delta\mu/e$——伏特是电子化学的
兑换率，这也是电池只能提供个位数伏特的原因。
\end{definition}

\section{一次一个能级：主技巧}

\begin{theorem}[单个量子能级的占据数]\label{thm:b3:grand-canonical:occupation}
\index{费米–狄拉克分布}\index{玻色–爱因斯坦分布}
取“体系”为能量为
$\epsilon$ 的一个单粒子能级，与其余所有能级构成的气体处于巨正则接触。
\emph{费米子}（占据数为 $0$ 或 $1$）：
\[
\Xi = 1 + \eu^{-\beta(\epsilon - \mu)}
\;\Longrightarrow\;
\bar n_{\text{FD}} = \frac{1}{\eu^{\beta(\epsilon - \mu)} + 1} .
\]
\emph{玻色子}（占据数为 $0, 1, 2, \dots$）：
\[
\Xi = \frac{1}{1 - \eu^{-\beta(\epsilon - \mu)}}
\;\Longrightarrow\;
\bar n_{\text{BE}} = \frac{1}{\eu^{\beta(\epsilon - \mu)} - 1} ,
\]
要求 $\mu < \epsilon_{\min}$。二者在占据数很小时都
退化为玻尔兹曼因子 $\eu^{-\beta(\epsilon - \mu)}$。
这两个公式各用两行即可导出，是量子统计的全部
输入；下一章正是对它们的收获。
\end{theorem}

\begin{proof}
费米子：两项。玻色子：等比级数，仅当
$\mu < \epsilon$ 时收敛。两种情形下 $\bar n = k_{\text{B}}T\,
\partial\ln\Xi/\partial\mu$ 即得所述形式。
\end{proof}

\begin{omfigure}
\begin{tikzpicture}
  \begin{axis}[omaxis, axis x line=bottom, axis y line=left,
    width=11.2cm, height=5.6cm, xmin=-6, xmax=6, ymin=0, ymax=2.3,
    xlabel={$(\epsilon - \mu)/k_{\text{B}}T$}, ylabel={$\bar n$},
    xlabel style={at={(axis description cs:0.5,-0.1)}, anchor=north},
    xtick={-4,-2,0,2,4}, ytick={0.5,1,2}, clip=false,
    legend style={font=\scriptsize, at={(0.97,0.93)}, anchor=north east,
      draw=none, fill=none}]
    \addplot[omDef, very thick, domain=-6:6, samples=200]
      {1/(exp(x)+1)};
    \addlegendentry{费米–狄拉克}
    \addplot[omThm, very thick, domain=0.45:6, samples=200]
      {1/(exp(x)-1)};
    \addlegendentry{玻色–爱因斯坦}
    \addplot[omExo, dashed, thick, domain=-0.8:6, samples=100]
      {exp(-x)};
    \addlegendentry{玻尔兹曼}
    \draw[gray, dashed] (axis cs:0,0) -- (axis cs:0,2.3);
  \end{axis}
\end{tikzpicture}

{\small 三种占据律关于 $(\epsilon -
\mu)/k_{\text{B}}T$ 的图像：费米子每个态至多占据一个（
在 $\epsilon = \mu$ 处为平滑阶跃），玻色子无限堆积，
当 $\epsilon \to \mu$ 时，二者在右侧都汇入稀薄的玻尔兹曼尾部——
经典区域正是一个态很少被占据的地方，
几乎从不被占据。}
\end{omfigure}

\section{吸附：租用的位点}

\begin{proposition}[朗缪尔等温线]\label{prop:b3:grand-canonical:langmuir}
\index{朗缪尔等温线}\index{吸附}
表面提供相互独立的位点，每个位点或空着，或\omterm{prop:b3:grand-canonical:langmuir}{吸附}一个气体
分子，其结合能为 $-\epsilon_0$；上方的气体
压强为 $P$，是能量库。每个位点是一个两态巨正则
体系，其占据数为
\[
\theta = \frac{1}{\eu^{\beta(-\epsilon_0 - \mu)} + 1}
 = \frac{P}{P + P_0(T)} , \qquad
P_0(T) = \frac{k_{\text{B}}T}{\lambda_T^3}\,
\eu^{-\epsilon_0/k_{\text{B}}T} ,
\]
使用理想气体的 $\mu$。低压下覆盖度线性增长，
并在单分子层处饱和——即\emph{\omterm{prop:b3:grand-canonical:langmuir}{朗缪尔等温线}}，
它是催化转化器、活性炭、防毒面具以及血液携氧
过程所遵循的工作曲线
（见\cref{exo:b3:grand-canonical:6}）。过渡压强 $P_0$
随结合能\emph{指数式}下降，并
随温度急剧上升：加热表面会把其上的“住客”赶走——
这正是真空烘烤玻璃器皿能真正清洁它的原因。
\end{proposition}

\begin{proof}
位点的占据数为
\cref{thm:b3:grand-canonical:occupation}（费米子形式：每个位点
至多容纳一个粒子），取 $\epsilon = -\epsilon_0$；代入
$\eu^{\beta\mu} = n\lambda_T^3 = P\lambda_T^3/k_{\text{B}}T$.
\end{proof}

\begin{omfigure}
\begin{tikzpicture}
  \begin{axis}[omaxis, axis x line=bottom, axis y line=left,
    width=10.8cm, height=5.4cm, xmin=0, xmax=5, ymin=0, ymax=1.12,
    xlabel={低温时的 $P/P_0$}, ylabel={覆盖度 $\theta$},
    xlabel style={at={(axis description cs:0.5,-0.1)}, anchor=north},
    xtick={1,2,3,4}, ytick={0.5,1}, clip=false,
    legend style={font=\scriptsize, at={(0.97,0.42)}, anchor=north east,
      draw=none, fill=none}]
    \addplot[omDef, very thick, domain=0:5, samples=120] {x/(x+1)};
    \addlegendentry{低温表面}
    \addplot[omThm, very thick, domain=0:5, samples=120] {x/(x+6)};
    \addlegendentry{高温表面（$P_0$ 更大）}
    \draw[gray, dashed] (axis cs:0,1) -- (axis cs:5,1);
    \node[font=\scriptsize, anchor=south east, gray] at (axis cs:4.9,1.0)
      {单分子层};
  \end{axis}
\end{tikzpicture}

{\small \omterm{prop:b3:grand-canonical:langmuir}{朗缪尔等温线}：覆盖度随压强变化，
在单分子层处饱和。加热表面后，要达到相同的覆盖度
需要高得多的压强——\omterm{prop:b3:grand-canonical:langmuir}{吸附}是一场 $\mu$ 的协商，而高温表面
总是输掉这场协商。}
\end{omfigure}

\section{电离平衡：萨哈方程}

\begin{proposition}[萨哈方程]\label{prop:b3:grand-canonical:saha}
\index{萨哈方程}
在热气体中，反应 $\text{H} \rightleftharpoons \text{p} +
\text{e}$ 达到平衡的条件为 $\mu_{\text{H}} = \mu_{\text{p}} +
\mu_{\text{e}}$。用理想气体形式写出每个 $\mu$（
原子计入其结合能 $-E_{\text{I}}$），得到
\[
\frac{n_{\text{e}}\,n_{\text{p}}}{n_{\text{H}}}
 = \frac{1}{\lambda_{T,\text{e}}^3}\,
\eu^{-E_{\text{I}}/k_{\text{B}}T}
\qquad\Big(\lambda_{T,\text{e}} =
h/\sqrt{2\pi m_{\text{e}}k_{\text{B}}T}\Big)
\]
（至多相差数量级为一的自旋简并因子）。这里又有两种相互竞争的“货币”：
玻尔兹曼因子对电离征收
$E_{\text{I}} = \qty{13.6}{eV}$ 的“税”，而前因子——
\emph{被释放电子的熵}，连同其 $\lambda_T^{-3}$
所对应的可及态数——则给予电离“补贴”。这笔补贴
在恒星密度下大得足以使氢
在 \qty{10000}{K} 附近就已显著电离，此时 $k_{\text{B}}T$ 仅为
\qty{0.9}{eV}：平衡之争在于\omterm{prop:b3:canonical-ensemble:machine}{自由能}，而非能量——
而恒星的光谱正是这本账本
（见\cref{pb:b3:grand-canonical:1}）。
\end{proposition}

\begin{proof}[部分证明]
令 $\mu_{\text{H}} = \mu_{\text{p}} + \mu_{\text{e}}$，其中 $\mu_i
= k_{\text{B}}T\ln(n_i\lambda_{T,i}^3) + \varepsilon_i$
（$\varepsilon_{\text{H}} = -E_{\text{I}}$，其余为零）；
两边取指数。质子与原子的热波长近似相消
（质量相等）；简并因子在本阶段不加证明地承认。
\end{proof}

\begin{omfigure}
\begin{tikzpicture}
  \begin{axis}[omaxis, axis x line=bottom, axis y line=left,
    width=10.8cm, height=5.4cm, xmin=3000, xmax=25000, ymin=0, ymax=1.1,
    xlabel={$T$ (K)}, ylabel={氢的电离比例},
    xlabel style={at={(axis description cs:0.5,-0.1)}, anchor=north},
    xtick={5000,10000,15000,20000}, ytick={0.5,1}, clip=false]
    \addplot[omDef, very thick, domain=3000:25000, samples=300]
      {1/(1+exp(-(x-9600)/1100))};
    \draw[gray, dashed] (axis cs:9600,0) -- (axis cs:9600,0.5);
    \node[font=\scriptsize, anchor=west] at (axis cs:10000,0.28)
      {\qty{10000}{K} 附近半电离};
  \end{axis}
\end{tikzpicture}

{\small 氢在光球层密度下的萨哈电离（示意）：
在 \qty{10000}{K} 附近出现急剧转变，远低于朴素估计
$E_{\text{I}}/k_{\text{B}} \approx \qty{160000}{K}$——被释放电子的
熵承担了大部分代价。}
\end{omfigure}

\begin{method}[巨正则技巧]\label{met:b3:grand-canonical:method}
(1) 凡与大系统交换粒子者：赋予其
$\mu$，并以 $\eu^{-\beta(E - \mu N)}$ 对状态加权。
(2) 独立的格点或能级可因子分解：求解\emph{一个}，
再相乘。(3) 对经典气体库，$\eu^{\beta\mu} =
n\lambda_T^3$ 将 $\mu$ 换算为可测压强。(4) 反应
与相变：平衡各物种的 $\mu$（每种一个），
计入结合能——预期熵的前因子会使平衡
远离朴素的玻尔兹曼估计。(5) 固体中的电子：
其 $\mu$ 即费米能级，$\mu$ 之差
即电压——与
\cref{ch:b3:electrons-in-solids}相衔接。
\end{method}

\section{练习}

\begin{exercise}[$\star$]\label{exo:b3:grand-canonical:1}
(a) 两个盛有同种气体的容器，温度同为 $T$ 而密度不同，相互连通：
粒子朝哪个方向流动？
用 $\mu$ 语言表述。(b) 由理想气体的 $\mu$ 证明流动
在密度相等时停止。(c) 水置于潮湿空气之下，
请以 $\mu$ 写出平衡条件。(d) 为何在空气干燥时，
衣物在低于 $100\,^\circ$C 时仍能晾干（比较 $\mu$，
而非温度）？
\end{exercise}

\begin{exercise}[$\star$]\label{exo:b3:grand-canonical:2}
空气的\omterm{thm:b3:grand-canonical:distribution}{化学势}。(a) 计算 $\lambda_T$（对象为 N$_2$）
在 \qty{300}{K} 下的值。(b) 计算 \qty{1}{bar} 下的 $n\lambda_T^3$，
并以 eV 表示所得的 $\mu$。(c) 为何 $\mu < 0$ 并不矛盾
（加入一个分子时，$F = E - TS$ 中哪一项占主导？）(d)
在何种密度下 $\mu$ 恰为零，
以及 $n\lambda_T^3 \sim 1$ 预示什么
（见\cref{exo:b3:grand-canonical:11}）？
\end{exercise}

\begin{exercise}[$\star$]\label{exo:b3:grand-canonical:3}
由其两项式的 $\Xi$ 逐步推导费米–狄拉克占据数；
检验极限 $\epsilon \ll \mu$，$\epsilon \gg
\mu$，以及 $\epsilon = \mu$ 处的取值。此函数
在本书中何处已经出现过
（见\cref{exo:b3:microcanonical-ensemble:8}）？
\end{exercise}

\begin{exercise}[$\star$]\label{exo:b3:grand-canonical:4}
电压与化学有关。(a) 一节锂电池的电压为 \qty{3.0}{V}：
请以 eV 表示每个电子的\omterm{thm:b3:grand-canonical:distribution}{化学势}差，并
以 kJ 表示每摩尔的\omterm{thm:b3:grand-canonical:distribution}{化学势}差。(b) 为何日常电池的电压都只有个位数
伏特（是什么决定了化学 $\mu$ 之差的量级？）(c) 一块铅酸
汽车电池由六节单格串联而成：为何要串联，而不直接
寻找一个 \qty{12}{eV} 的反应？(d) 以 $\mu$ 语言描述，
电池在“没电”时正在做什么？
\end{exercise}

\begin{exercise}[$\star\star$]\label{exo:b3:grand-canonical:5}
一个能级上的玻色子。(a) 对 $\Xi$ 的几何级数求和，并
推导 $\bar n_{\text{BE}}$。(b) 为何 $\mu$ 必须低于最低能级
——若不然，什么会发散，发散又
预示了什么（见\cref{ch:b3:quantum-statistics}）？(c) 计算
$\bar n$，取 $(\epsilon - \mu)/k_{\text{B}}T = 0.1$、$1$、
$3$ 的情形，并与玻尔兹曼结果比较。(d) 光子的 $\mu = 0$，请给出
物理原因（光子数不守恒——器壁
可自由地吸收与发射光子）。
\end{exercise}

\begin{exercise}[$\star\star$]\label{exo:b3:grand-canonical:6}
朗缪尔的应用。(a) 由等温线，
使 $\theta = 0.5$、$0.9$ 的压强分别为多少？(b) 肌红蛋白在一个位点上结合 O$_2$：
饱和度随氧分压的变化遵循朗缪尔规律——
为何肌红蛋白呈现的是\emph{双曲}（而非 S 形）曲线，
而四位点的血红蛋白却呈协同的 S 形曲线？
(c) 一氧化碳结合于相同位点，$\epsilon_0$
高出 $\approx\qty{0.13}{eV}$，在分压相等
及 \qty{310}{K} 下，计算占据比 CO:O$_2$，并解释
一氧化碳中毒。(d) 为何高压氧疗有效，
用等温线语言说明？
\end{exercise}

\begin{exercise}[$\star\star$]\label{exo:b3:grand-canonical:7}
粒子数涨落。(a) 证明 $\operatorname{Var}N =
k_{\text{B}}T\,\partial\langle N\rangle/\partial\mu$。(b) 对
经典理想气体，推出 $\operatorname{Var}N = \langle
N\rangle$，即泊松分布。(c) 使其与
\cref{exo:b3:microcanonical-ensemble:12}的密度涨落相协调。(d) 对单个
费米子能级，计算 $\operatorname{Var}n = \bar n(1 - \bar
n)$：费米子数在何处涨落最小，原因何在？
\end{exercise}

\begin{exercise}[$\star\star$]\label{exo:b3:grand-canonical:8}
由 $\mu$ 理解渗透。一层膜允许水透过、阻挡溶质
（单位体积浓度为 $c$，稀溶液）。(a) 令膜两侧水的 $\mu$ 相等，
证明纯水一侧须多承受
超压 $\Pi = ck_{\text{B}}T$（范特霍夫）——并接受溶质会使
水的 $\mu$ 降低
$-k_{\text{B}}T\,c/n_{\text{w}}$（每个水分子）。(b) 计算
生理盐水的 $\Pi$（$c \approx \qty{3e26}{m^{-3}}$
总粒子数）。(c) 红细胞置于纯水中：水向何方流动，
产生何种结果？(d) 树木为何既不需要
泵，也无需奇迹，即可将汁液提升数十米（渗透与
毛细作用造成的 $\mu$ 梯度）？
\end{exercise}

\begin{exercise}[$\star\star$]\label{exo:b3:grand-canonical:9}
场中的平衡。在重力场中，气柱的平衡
要求\emph{总}\omterm{thm:b3:grand-canonical:distribution}{化学势} $\mu(h) =
k_{\text{B}}T\ln(n\lambda_T^3) + mgh$ 须为常数。(a) 由此在一行内推出
气压随高度的分布律。(b) 推广至离心机，
设角速度为 $\omega$（势能
$-\tfrac12 m\omega^2r^2$），并求径向分布。(c) 铀浓缩
离心机的线速度约为 $\omega r \sim
\qty{600}{m/s}$，计算平衡时 $^{238}$UF$_6$
与 $^{235}$UF$_6$ 之比的每级增强因子，$\exp[\Delta m\,\omega^2
r^2/2k_{\text{B}}T]$，取 $\Delta m = \qty{5e-27}{kg}$，$T =
\qty{320}{K}$。(d) 为何要采用多级级联？
\end{exercise}

\begin{exercise}[$\star\star\star$]\label{exo:b3:grand-canonical:10}
手算\omterm{prop:b3:grand-canonical:saha}{萨哈方程}。(a) 完成
\cref{prop:b3:grand-canonical:saha}的 $\mu$ 平衡证明。(b) 定义电离度 $x$
（总氢数密度为 $n$）并证明 $x^2/(1 - x)
= \eu^{-E_{\text{I}}/k_{\text{B}}T}/n\lambda_{T,\text{e}}^3$。(c)
太阳光球层：$T = \qty{5800}{K}$，$n \approx
\qty{e23}{m^{-3}}$，证明氢的电离度仅为
$\sim\num{e-4}$——太阳表面呈中性。(d) 在 $T =
\qty{10000}{K}$（$n$ 不变）：重新计算，并评论
其陡峭程度。
\end{exercise}

\begin{exercise}[$\star\star\star$]\label{exo:b3:grand-canonical:11}
气体何时进入量子区。简并出现于 $n\lambda_T^3 \sim
1$。(a) 对于 \qty{300}{K} 的空气，低于阈值多少
（见\cref{exo:b3:grand-canonical:2}）？(b) 对于铜中的导电电子
（$n = \qty{8.5e28}{m^{-3}}$）：求温度
低于该温度时 $n\lambda_T^3 > 1$——并由此得出金属电子在
\emph{任何}实验室温度下都是简并的。
(c) 对于 $n = \qty{e20}{m^{-3}}$ 的铷原子，求量子
阈值温度（即
\cref{ch:b3:quantum-statistics}中的玻色–爱因斯坦凝聚特征温度）。(d) 将三者排序，并陈述规律：
轻粒子和高密度最先进入量子区。
\end{exercise}

\begin{exercise}[$\star\star\star$]\label{exo:b3:grand-canonical:12}
宇宙失去了它的正电子。在炽热的早期宇宙中，
$\eu^+ + \eu^- \rightleftharpoons 2\gamma$ 使正负电子对维持平衡，
平衡条件为 $\mu_{e^+} + \mu_{e^-} = 0$（光子：$\mu = 0$）。
(a) 对于对称等离子体，由此得 $\mu_{e^\pm} \approx 0$，
并证明此时正负电子数密度为 $n_\pm \approx
\lambda_T^{-3}\eu^{-m_{\text{e}}c^2/k_{\text{B}}T}$，当
$k_{\text{B}}T \ll m_{\text{e}}c^2$ 时。(b) 求温度，
使 $m_{\text{e}}c^2 = k_{\text{B}}T$，并给出它所对应的宇宙时间尺度
（可取 $t \sim \qty{1}{s}$ 对应
$\qty{e10}{K}$）。(c) 证明进一步适度降温（降幅约
几倍）即可使正负电子对数密度
下降许多数量级：正电子湮灭殆尽，只留下原初的
微小电子过剩。(d) 释放出的光子再加热
辐射浴：哪一种遗留辐射保存了这一
记录（见\cref{ch:b3:astrophysics}）？
\end{exercise}

\section{问题：读取恒星的温度}

\begin{problem}[{周末问题——萨哈、佩恩与
宇宙由何构成}]\label{pb:b3:grand-canonical:1}
1890 年代，恒星光谱被编目为
O、B、A、F、G、K、M 等类，依据其吸收线的强弱——
其中氢的巴耳末线在 A 类中最强，在较热与较冷的
恒星中却都神秘地偏弱。1925 年，塞西莉亚·佩恩
运用刚提出的\omterm{prop:b3:grand-canonical:saha}{萨哈方程}，证明这些光谱类型
其实是同一种物质在不同温度下的表现——并且恒星
绝大部分由氢构成，这一结论如此离经叛道，以至于她被迫
在学位论文中予以淡化。本题重现她的推理过程。
数据：巴耳末吸收始于 $n = 2$，能级差 $E_2 - E_1 =
\qty{10.2}{eV}$（高于基态）；$E_{\text{I}} =
\qty{13.6}{eV}$；光球层电子数密度 $n_{\text{e}}
\sim \qty{e20}{m^{-3}}$。

\medskip\noindent\textbf{第一部分——两个相互竞争的因素。}
\begin{enumerate}
  \item 巴耳末吸收需要中性氢\emph{处于} $n =
        2$）：写出 $n = 2$ 能级布居的玻尔兹曼因子
        （简并度：$g_2/g_1 = 4$）。
  \item 在 $5000$、$10000$ 与 \qty{20000}{K} 处计算该因子：
        该因子随 $T$ 升高如何改变谱线强度？
  \item 再看萨哈因子：随 $T$ 升高，
        \emph{中性}比例会如何变化，进而巴耳末线又如何？
  \item 定性解释为何这一乘积——即既被激发
        \emph{又}仍为中性——必在某个
        中间温度处达到峰值。
  \item 利用\omterm{prop:b3:grand-canonical:saha}{萨哈方程}，取 $n_{\text{e}} =
        \qty{e20}{m^{-3}}$，估算氢
        半电离温度。
  \item 综合上述结果：巴耳末线应在何温度附近最强，
        以及这一温度挑出的是哪一光谱类型
        （A 型星，$\sim\qty{10000}{K}$）？
\end{enumerate}

\medskip\noindent\textbf{第二部分——光谱序列的
解读。}
\begin{enumerate}[resume]
  \item M 型星（\qty{3000}{K}）呈现强分子带（TiO）
        并且氢线偏弱：请用这两个因素解释这两点。
  \item O 型星（\qty{40000}{K}）呈现氦线，且氢线同样偏弱：
        请解释（氢发生了什么，以及为何
        氦（$E_{\text{I}} = \qty{24.6}{eV}$）此时
        轮到它出场？）
  \item 太阳（G 型，\qty{5800}{K}）呈现强烈的
        \emph{电离钙}谱线（$E_{\text{I,Ca}} = \qty{6.1}{eV}$）
        而氢线偏弱：请用萨哈逻辑说明为何钙已
        电离，而氢仍为中性，且主要位于
        $n = 1$ 能级。
  \item 佩恩之前的天文学家把谱线强度当作丰度：
        太阳强大的钙线与微弱的巴耳末线
        使他们对其成分得出了怎样的结论？
  \item 佩恩的洞见：对每条谱线的强度，
        需修正其玻尔兹曼–萨哈“可见度因子”。她这样修正之后，
        真实丰度呈现出怎样的层级？
  \item 她的结论——氢的丰度是钙的百万倍，
        恒星主要由氢和氦构成——为何它遭到
        抵制（地心化学暗示了什么），以及她何时
        得到证实？
\end{enumerate}

\medskip\noindent\textbf{第三部分——峰值上的数值。}
\begin{enumerate}[resume]
  \item 计算 $n = 2$ 的玻尔兹曼因子，温度恰为
        \qty{9600}{K}。
  \item 利用第 5 项的萨哈估算，取
        $9600$、$15000$ 与 \qty{25000}{K} 下的中性比例近似为
        $0.5$、$\num{e-3}$、$\num{e-5}$：
        计算三个温度下的乘积（激发比例 $\times$ 中性比例）
        并确定其极大值位置。
  \item 为何巴耳末峰在高温一侧\emph{陡峭}
        （哪一项指数占优），而在低温一侧则
        \emph{平缓}？
  \item 白矮星光球层的 $n_{\text{e}}$ 大一千倍：
        半电离温度会向哪个方向移动，
        以及为什么（以 $\mu$ 语言表述的
        勒夏特列原理）？
  \item 两颗温度相同、压强不同的恒星
        因此呈现出不同的谱线强度：
        这使光谱学家能读出哪一项恒星性质（巨星与矮星之别
——即表面重力）？
  \item 现代巡天依据谱线比值，将恒星温度确定到 $\pm1\%$，
        请陈述推理链：假设 $\to$ 玻尔兹曼
        $\to$ 萨哈 $\to$ 温度计。
\end{enumerate}

\medskip\noindent\textbf{第四部分——所确立的内容。}
\begin{enumerate}[resume]
  \item OBAFGKM 序列原本是字母顺序的混乱排列，
        经验重排之后：在萨哈之后，哪一个单一的物理变量
        决定了它的排序？
  \item 说明为何\emph{同一}元素可在某一温度下主导光谱，
        而在另一温度下却消失，丰度却没有任何变化
——这正是佩恩教给天文学界的
        核心道理。
  \item 氢占全部重子质量的四分之三：试指出本书另外两个
        建立在这一丰度之上的支柱
        （恒星聚变，\cref{ch:b3:astrophysics}；21 厘米天空，
        \cref{ch:b3:spin-two-level}）。
  \item 奥托·斯特鲁维称佩恩的博士论文为
        “天文学史上写得最出色的博士论文”：
        请依据本题，用一句话论证这一赞誉。
  \item 同样的萨哈物理，若在宇宙尺度上反向运用，
        可预言整个宇宙的氢首次复合的温度
        （宇宙密度远低于此，约为 $\sim\qty{3000}{K}$）：
        指出该事件及其遗迹
        （\cref{ch:b3:astrophysics}）。
  \item 电离平衡同样支配荧光灯与等离子体加工：
        请用一句话说明，为何一根玻璃管保持低温，
        管内却能存在类似 \qty{10000}{K} 的电离态，
        （哪两个子系统未能达到平衡
        ？）。
  \item 概括这一命名结果：一个平衡条件，
        $\mu_{\text{H}} = \mu_{\text{p}} + \mu_{\text{e}}$，使
        光谱字母表化为一支温度计，量程从
        $3000$ 到 \qty{40000}{K}，使巴尔末峰落在 A 型星上，
        并揭示宇宙由氢构成。
\end{enumerate}
\end{problem}
